%% ma: L11-B16-0068
%% lop: 11
%% chuong: 5
%% bai: 16
%% dang: 11.16.2
%% loai: TLN
%% muc: VD
%% dap_an: "16"
%% nguon: "tai-lieu-hoc-tap-mon-toan-lop-11-theo-chuong-trinh-giao-duc-pho-thong-2018.pdf (Lớp Toán Thầy Duy, Phạm Lê Duy, THPT Chu Văn An), Chương 3 Bài 2, Phần 3, Câu 6 (tr.381)"
%% kien_thuc: "Lớp 11 Bài 16 (giới hạn của hàm số)"
%% trang_thai: chinh-thuc
%% ghi_chu: "Đã duyệt 03/11/2026."
\begin{ex}
Biết rằng giới hạn \[\lim\limits_{x\to2}\dfrac{\sqrt{5x-1}-\sqrt{9x-2}+1}{x-2}=\lim\limits_{x\to2}\left(\dfrac a{b\sqrt{5x-1}+3}-\dfrac c{4+d\sqrt{9x-2}}\right)\] với $a,b,c,d$ là các số nguyên dương. Tính tổng các số $a,b,c,d$.
\shortans{16}
\loigiai{\sloppy Vì $3-4+1=0$ nên \[\sqrt{5x-1}-\sqrt{9x-2}+1=\left(\sqrt{5x-1}-3\right)-\left(\sqrt{9x-2}-4\right)=\dfrac{5(x-2)}{\sqrt{5x-1}+3}-\dfrac{9(x-2)}{\sqrt{9x-2}+4}.\] Chia cho $x-2$: biểu thức $=\dfrac5{\sqrt{5x-1}+3}-\dfrac9{4+\sqrt{9x-2}}$. Vậy $a=5$, $b=1$, $c=9$, $d=1$ và tổng $=16$.}
\end{ex}
